
Sabemos que
\begin{equation*}
    L_t(\hat{\alpha}_t, \hat{\theta}_t) = 2\sqrt{W_t^= W_t^{\neq}} \le L_t(\tilde{\alpha}, \hat{\theta}),
\end{equation*}
onde
\begin{equation*}
    L_t(\tilde{\alpha}, \hat{\theta}) = \sum_{i=1}^n w_{it} \, e^{-\tilde{y}_i \tilde{\alpha} F(x_i, \hat{\theta})}.
\end{equation*}

Separando os termos de acordo com $\tilde{y}_i = F(x_i, \hat{\theta})$ e $\tilde{y}_i \neq F(x_i, \hat{\theta})$, obtemos
\begin{equation*}
    L_t(\tilde{\alpha}, \hat{\theta}) = W_t^= \left( \frac{1-2\gamma}{1+2\gamma} \right)^{1/2} + W_t^{\neq} \left( \frac{1+2\gamma}{1-2\gamma} \right)^{1/2}.
\end{equation*}

Fatorando $\sqrt{1-4\gamma^2}$, temos
\begin{equation}
    L_t(\tilde{\alpha}, \hat{\theta}) = \sqrt{1-4\gamma^2} \left( \frac{W_t^=}{1+2\gamma} + \frac{W_t^{\neq}}{1-2\gamma} \right) \tag{I}
\end{equation}

Do item 2d, sabemos que
\begin{equation*}
    W_t^{\neq} \le \frac{1-2\gamma}{1+2\gamma} W_t^=.
\end{equation*}

Substituindo essa desigualdade em \eqref{I}, segue que
\begin{equation*}
    \frac{W_t^{\neq}}{1-2\gamma} \le \frac{W_t^=}{1+2\gamma}.
\end{equation*}

Portanto,
\begin{equation*}
    \frac{W_t^=}{1+2\gamma} + \frac{W_t^{\neq}}{1-2\gamma} \le \frac{2 W_t^=}{1+2\gamma}.
\end{equation*}

Como
\begin{equation*}
    W_t^= + W_t^{\neq} \le \frac{2 W_t^=}{1+2\gamma},
\end{equation*}
concluímos que
\begin{equation*}
    L_t(\hat{\alpha}_t, \hat{\theta}_t) \le \sqrt{1-4\gamma^2} \left( W_t^= + W_t^{\neq} \right) = \sqrt{1-4\gamma^2} \, L_{t-1}(\hat{\alpha}_{t-1}, \hat{\theta}_{t-1}).
\end{equation*}

Inicializando $\hat{\alpha}_0 = 0$, de modo que $L_0(\hat{\alpha}_0, \hat{\theta}_0) = 1$, e usando a desigualdade $1+x \le e^x$ para todo $x \in \mathbb{R}$, obtemos
\begin{equation*}
    L_t(\hat{\alpha}_t, \hat{\theta}_t) \le e^{-2 t \gamma^2}.
\end{equation*}

Assim, o erro de treino do AdaBoost com a perda exponencial decresce de forma exponencialmente rápida.